\section{Nuevas reflexiones sobre el scaling}

\subsection{Cómo el lenguaje moldea el pensamiento}

Ilya planteó una observación profunda: \textbf{el lenguaje influye en el pensamiento}. Basta la palabra «Scaling» para impulsar el comportamiento de toda la industria: la gente dice «let's scale» y luego todos hacen lo mismo.

\begin{importantbox}{El doble legado del Pre-training}
Dos términos moldearon profundamente el pensamiento de todo el campo:
\begin{enumerate}
    \item \textbf{AGI}: esta palabra surgió como reacción a la «narrow AI» (IA de ajedrez, IA de videojuegos); dado que la IA estrecha no era suficiente, se necesitaba una «general AI»
    \item \textbf{Pre-training}: su recipe era «more pre-training $\rightarrow$ better at everything, uniformly»
\end{enumerate}
La combinación de ambos produjo la mentalidad fija de «pre-training gives AGI». Pero el problema es: \textbf{los propios humanos no son AGI}; a los humanos les falta una gran cantidad de conocimiento y dependen del continual learning.
\end{importantbox}

\subsection{La era de Research no necesita el cómputo a máxima escala}

\begin{knowledgebox}{El cómputo necesario para los avances clave en la historia}
\begin{itemize}
    \item AlexNet: 2 GPU
    \item El artículo de Transformer: como máximo 64 GPU (de 2017, equivalente aproximadamente a 2 GPU actuales)
    \item ResNet: una escala pequeña similar
    \item O1 style reasoning: tampoco es lo más intensivo en cómputo del mundo
\end{itemize}
En la etapa de investigación, no necesitas la escala de cómputo absolutamente más grande para demostrar que una idea es correcta. El cómputo a gran escala es más bien un factor de diferenciación cuando todos compiten dentro del mismo paradigm.
\end{knowledgebox}

\subsection{El posicionamiento de cómputo de SSI}

Ilya explicó por qué el financiamiento de \$3,000 millones de SSI no es poco para la investigación:
\begin{itemize}
    \item La mayor parte del cómputo de las grandes empresas se destina a inference (para dar servicio a productos)
    \item Las grandes empresas necesitan una gran cantidad de personal (ingenieros, ventas) e investigación relacionada con productos
    \item La cantidad de cómputo realmente destinada a la investigación pura tiene una brecha mucho menor que la que sugieren las cifras superficiales
\end{itemize}

\subsection{Resumen de la sección}

La era de Scaling «se llevó todo el aire de la habitación», lo que provocó que hubiera más empresas que ideas. Volver a la era de research significa volver a valorar la creatividad y la experimentación, no solo la escala. Los avances clave no necesariamente requieren el cómputo a la máxima escala.


